% Chapter converted from Markdown
% Auto-generated - do not edit manually

\# Chapter 2 - Vision: Chat/IDE as the Universal Interface


\subsection*{Executive Summary}

\item The interface decides whether AI work is improvable. Chat must become the control plane; the IDE remains the substrate where artifacts, tests, and evidence live.
\item A universal Chat/IDE surface wires together the core systems introduced in Chapter 1: CMC, HHNI, VIF, APOE, and SEG.
\item The vision is meta-circular: we use the interface to write and validate this chapter, demonstrating that the tools already exist and that the workflow is reproducible.

\subsection*{Interface Principles}

| Principle | Description | Resulting Behavior |
| --- | --- | --- |
| \textbf{Statefulness} | Chat threads resume with intent, plans, and atoms retrieved from CMC. | Work restarts with context loaded instead of manual recap. |
| \textbf{Constraint-first} | Gates and policies shape the conversation. Plans, evidence, and examples are negotiated explicitly. | Quality is enforced in-line, not as a review afterthought. |
| \textbf{Shared visibility} | Agents and humans see the same files, metrics, and tool outputs inside the IDE. | Collaboration becomes computable and auditable. |
| \textbf{Runnable truth} | Every major claim pairs with a runnable example or validated script. | The system proves its capability as it describes it. |

\subsection*{Roles of Chat and IDE}

\item \textbf{Chat stream:} sets objectives, proposes plans, records status, and routes confidence updates. Messages link directly to artifact diffs, tests, and evidence entries.
\item \textbf{IDE workspace:} stores chapters, code, metrics, and evidence. MCP tools operate on the workspace with policy-enforced safety checks.
\item \textbf{Command surfaces:} operations like \texttt{run\_autonomous\_checklist} and \texttt{get\_tag\_coverage} expose reproducible controls. Results are written to files plus surfaced in chat summaries.

\subsection*{Capabilities the Interface Must Provide}

1. \textbf{Memory Access (CMC):} Retrieve atoms mapped to the active goal, scope, and time horizon without manual searching.
2. \textbf{Hierarchical Retrieval (HHNI):} Navigate from executive summaries to deep-dive nodes in a few jumps, maintaining coherence.
3. \textbf{Confidence Routing (VIF):} Record and enforce thresholds. If confidence drops below policy, the system stops or diverts to research automatically.
4. \textbf{Executable Plans (APOE):} Produce chains that specify steps, expected artifacts, validation hooks, and escalation logic.
5. \textbf{Evidence Graph (SEG):} Attach claims to anchors, detect contradictions, and block merges when proof is missing.

\subsection*{Runnable Interface Examples (PowerShell)}

Discover relevant tools for the current workspace:

``\texttt{powershell
Invoke-WebRequest -Uri 'http://localhost:5001/mcp/list' -Method GET |
  Select-Object -ExpandProperty Content | ConvertFrom-Json |
  ForEach-Object { \_\_MATH\_BLOCK\_0\_\_stats = @{ tool='get\_memory\_stats'; arguments=@{} } | ConvertTo-Json -Depth 5
Invoke-WebRequest -Uri 'http://localhost:5001/mcp/execute' -Method POST -ContentType 'application/json' -Body \_\_MATH\_BLOCK\_1\_\_query = @{ tool='retrieve\_memory'; arguments=@{ query='universal interface vision'; limit=5 } } | ConvertTo-Json -Depth 6
Invoke-WebRequest -Uri 'http://localhost:5001/mcp/execute' -Method POST -ContentType 'application/json' -Body \_\_MATH\_BLOCK\_2\_\_uri = 'http://localhost:5001/mcp/execute'
\_\_MATH\_BLOCK\_3\_\_false
} } | ConvertTo-Json -Depth 6
Invoke-RestMethod -Uri \_\_MATH\_BLOCK\_4\_\_body | Out-Null
}`\texttt{

Example B — List project commands via MCP (discovery):

}`\texttt{powershell
\_\_MATH\_BLOCK\_5\_\_true } } | ConvertTo-Json -Depth 6
Invoke-RestMethod -Uri \_\_MATH\_BLOCK\_6\_\_body | ConvertTo-Json -Depth 5
}`\texttt{

Use the audited route (}/mcp/execute\texttt{) for MCP tools. Cite output identifiers in }evidence.jsonl\texttt{ when referencing results.

\subsection*{Operational Runbook (Minimal Loop)}

1) Check in (MCP), 2) edit + add one example, 3) append Tier A evidence, 4) run gates for this chapter, 5) post gate outcomes to the shared board. Small loops stay auditable and reversible.

\subsection*{Performance Characteristics (Local)}

\item Centralized: Command server handles }/mcp/execute\texttt{ and chat macros with logging.
\item Sub-second: Local MCP calls are typically quick; variance depends on environment.
\item Observable: Gate telemetry and server logs expose inputs and results.

\subsection*{Scenario: Coordinating Four Cursor Agents}

}north\_star\_project/CURSOR\_AGENT\_ONBOARDING.md\texttt{ shows how the same interface onboards Max, Lex, Sam, and Dac. Each agent sends the }send\_ai\_message\texttt{ payload from }AGENT\_CHECK\_IN\_PROTOCOL.md\texttt{, HHNI loads the relevant chains, and the IDE enforces runnable examples before posting to }coordination/epic\_standards\_overhaul/comms/SHARED\_MESSAGE\_BOARD.md\texttt{.

\subsection*{Intelligent Gate Telemetry}

Gate policy (}north\_star\_project/policy/gates.json\texttt{) replaced raw counts with relevance, density, completion, and thoroughness scores. Tier B chapters must keep relevance ≥0.82; missing examples trip the density gate and open a SIS task; completion stays }pending\texttt{ until Aether publishes the new spec. }north\_star\_project/scripts/run\_chain.py\texttt{ emits the same numbers that appear in }metrics.yaml\texttt{, so reviewers see the raw telemetry, not a guess.

\subsection*{Failure Modes and Mitigations}

\item \textbf{Tool overload:} limit surfaced tools to task-relevant capabilities via RAG filtering. Provide "why surfaced" explanations.
\item \textbf{Context drift:} HHNI enforces navigation discipline. Plans record chosen scope and depth.
\item \textbf{Gate fatigue:} automate checklists and run them opportunistically (on save, before merge).
\item \textbf{Evidence decay:} SEG monitors freshness and creates SIS tasks when anchors age beyond threshold.

\subsection*{Troubleshooting Guide}

\item 404 from command server: Confirm POST to }http://localhost:5001/mcp/execute\texttt{ and that the server is running.
\item Tool not visible: RAG filtering may hide it; provide clearer context or call }list\_cursor\_commands\texttt{.
\item Unicode/emoji crash: Use a UTF‑8 terminal or set }[Console]::OutputEncoding = [Text.UTF8Encoding]::new($false)\texttt{.
\item Cross-agent messaging not appearing: }send\_ai\_message` writes to both files; verify both JSON stores exist.

\subsection*{Demonstration: Minimal Edit Cycle}

1. Draft two sentences describing the change in chat.
2. Update the chapter file; add or adjust a runnable example.
3. Append an evidence entry with tier, source, and anchor.
4. Run contradiction and example gates; remediate issues immediately.
5. Post a status update summarizing the change plus confidence delta.

\subsection*{FAQ}

\textbf{Is this interface only for large teams?} No. The smallest loop is still tiny: a short objective, a runnable example, and one evidence entry. The ceremony scales down.

\textbf{Do contributors need to understand all subsystems?} The interface abstracts them. You only dive into details when troubleshooting or extending capability.

\textbf{Will gates slow down experts?} The discipline shortens review cycles and prevents rework. Time saved on regressions easily offsets gate execution.

\textbf{Can we customize policies?} Yes. Policy files define thresholds and escalation rules. Changes require evidence and review, preserving auditability.

\subsection*{Completeness Checklist (Chapter 2)}

\item Coverage complete: vision, principles, interaction loop, surfaces, scenarios, and FAQ.
\item Relevance sufficient: every section supports the claim that Chat/IDE must be the universal interface.
\item Subsection balance: no section dominates; conceptual and operational content share the space.
\item Minimum substance: runnable examples, tables, and timelines meet drafting requirements.

